\section{Conclusiones y trabajo futuro}
\label{sec:conclusiones}

La pregunta que abrió este artículo tenía una respuesta afirmativa, aunque
con matices. Un pipeline clásico, sin transformers ni representaciones
preentrenadas, puede competir en un corpus pequeño y ruidoso de tuits en
español si se cuidan tres cosas: la representación, la validación y la
decisión. El sistema alcanzó un $F_1$ público de 0.81447 y el primer lugar
del leaderboard al cierre de este informe. La estimación limpia de su
capacidad de generalización es, sin embargo, la del modelo base: 0.787582
out-of-fold, con una mejora de $+0.010370$ sobre V03 e IC95
$[+0.003924, +0.017386]$. Las dos cifras miden cosas distintas y ambas deben
reportarse.

La primera lección es que la representación sigue importando. La unión
multiescala de palabras y n-gramas de caracteres aportó $+0.00516$ local y
$+0.00549$ en el leaderboard, y lo hizo en las dos particiones congeladas.
En un corpus normalizado a ASCII, corto y con ruido ortográfico, los
n-gramas de caracteres no son un adorno: son el mecanismo que recupera la
morfología y la jerga que el vocabulario de palabras pierde.

La segunda lección es que la validación puede inducir a error si se diseña mal. P109
obtuvo la mejor métrica local de toda la serie y el peor resultado público.
La causa fue la poda por frecuencia: el vocabulario se hizo más pequeño,
la validación compartió esa estrechez y el test no. La validación agrupada,
la réplica en dos particiones y los controles de determinismo no eliminan el
problema, pero lo hacen visible. Sin ellos, P109 habría parecido el mejor
modelo del proyecto.

La tercera lección es que la capa de decisión tiene una economía propia. La
Proposición~\ref{prop:par} demuestra que un par compensado incrementa la
$F_1$ en $F_1/TP$ por columna y preserva el número de positivos. Esa fórmula
predice por qué los pares funcionan, por qué TR es la columna más sensible y
por qué P129 falló al saturar la probabilidad. También explica por qué la
capa no se puede validar con la OOF: sus reglas se diseñaron para textos del
test. El 76\,\% de la ganancia del leaderboard proviene de ahí, y esa es
precisamente la parte que exige más cautela.

El trabajo futuro tiene cuatro frentes. Primero, repetir el pipeline con
transformers o modelos de lenguaje, si las reglas lo permiten, para medir
cuánto de la brecha es representación y cuánto es arquitectura. Segundo,
atacar el error residual con mejor etiquetado y con modelos que manejen
ironía, discurso citado y odio implícito; el análisis por longitud muestra
que los textos largos son el frente principal. Tercero, diseñar una capa de
decisión que no dependa del leaderboard: calibración con validación
prospectiva, abstención cuando la probabilidad está cerca del umbral y
auditoría de subgrupos. Cuarto, evaluar el sistema con criterios de equidad
y robustez antes de pensar en un despliegue real; la literatura sobre sesgos
en corpus de odio es una advertencia que este trabajo no resuelve
\cite{davidson2019racial,wiegand2019biased}.

El resultado que queda es moderado pero significativo. Moderado porque
0.81447 es el puntaje de una competencia, con una capa adaptada al test y un
leaderboard privado que aún no se conoce. Significativo porque demuestra que
una representación construida desde cero, un ensamble lineal bien calibrado
y una capa de decisión auditable pueden llevar a un sistema clásico hasta el
primer lugar. En un momento en que la alternativa por defecto son los transformers, este
caso recuerda que la ingeniería de datos, la validación y
la teoría de la decisión siguen siendo fuentes reales de mejora.
